\section*{Codebeispiel für das Erstellen des Dropdownmenüs}
Der hier angezeigte Codeausschnitt erzeugt das Dropdownmenü auf unserer Hauptseite. Anschließend lässt sich über das ausgeklappte Menü ein Schiff auswählen. Wir haben uns bewusst dafür entschieden, dass ausschließlich der Name ausgewählt werden kann, da eine MMSI für die meisten nicht im Alltag vorkommt. Wenn man jedoch in oder um Bremerhaven wohnt, hat man bestimmt den ein oder anderen Schiffsnamen schon einmal gelesen.
\vspace*{0.5cm}
%listings als code - nicht als screenshots ...
\begin{lstlisting}[language=bash,caption={Schiffswelt.sh},captionpos=b,frame=tb]
#!/bin/bash
echo "Content-type: text/html"
echo
echo "<!doctype html><html><head><title>Schiffswelt</title><meta charset='utf-8'></head><body>"
cat /home/docker-step2022team12/rhodes.log | grep "^[^|]*|5|" | cut -d "|" -f 3,6 | sort -u -t "|" -k 2 > /home/docker-step2022team12/Schiffswelt/ships.csv
echo "<form action='shipdata.sh'><select name='Schiff'>"
i=1
cat /home/docker-step2022team12/Schiffswelt/ships.csv | while IFS="|" read mmsi name; do
echo "<option value=$i>$name</option>"
i=$((i+1))
done
echo "</select><input type=submit></form>"
echo "</body></html>"
\end{lstlisting}
